\begin{textsample}{POS Dim 1 – sp – Score 26.00 – sp/sp\_ef\_ciencias\_humanas\_geo\_3.txt}  \label{ex:f1_pos_169}
Diante do \textbf{exposto}, é \textbf{imprescindível} \textbf{que} o professor se reconheça como mediador no processo de \textbf{ensino-aprendizagem}, de forma \textbf{que} possa contribuir para a formação de cidadãos reflexivos, críticos, autônomos e transformadores da realidade local, regional e global, para a ampliação de repertório \textbf{teórico-metodológico} e para a formação integral dos estudantes.
Para \textbf{que} isso ocorra, é importante a apropriação de novos caminhos \textbf{metodológicos} para \textbf{um} processo de \textbf{ensino-aprendizagem} mais dinâmico, criativo e \textbf{interessante}.
Nos dias \textbf{atuais}, as metodologias ativas ( aprendizagem \textbf{baseada} em projetos, aprendizagem \textbf{baseada} em problemas, ensino híbrido, gamificação, entre outras ) são possibilidades para o fortalecimento do ensino de Geografia, \textbf{uma} vez \textbf{que} apresentam estratégias para o desenvolvimento das competências específicas do componente, da área de Ciências Humanas e de enfoques interdisciplinares e transversais.
Para o desenvolvimento dessas estratégias, é \textbf{imprescindível} \textbf{que} o professor busque aprimoramento constante da sua formação, de forma a consolidar a autonomia docente.
Ao mesmo tempo, é \textbf{preciso} \textbf{que} o estudante se reconheça como \textbf{um} \textbf{sujeito} \textbf{que} \textbf{vive} em \textbf{um} mundo contraditório e desafiador bem como suas responsabilidades na construção de \textbf{uma} sociedade \textbf{justa}, igualitária e sustentável.
Assim, os seus conhecimentos prévios, experiências, \textbf{percepções} e memórias individuais e coletivas são essenciais para a construção dos conhecimentos geográficos.
O desenvolvimento de conteúdos e temáticas relacionadas, por exemplo, à \textbf{crise} socioambiental, ao desenvolvimento econômico, às relações internacionais, à globalização, à diversidade cultural, aos desastres naturais, aos conflitos, ao agronegócio, às políticas públicas territoriais, às correntes migratórias, às mudanças climáticas, \textbf{aproximam} os estudantes de outras escalas de análise e fenômenos geográficos.
Assim sendo, ampliam o seu repertório de leitura de mundo e são estimulados a pensar espacialmente — tendo como referência os espaços cotidianos, espaços físicos e sociais — e a desenvolver os \textbf{raciocínios} geográficos \textbf{baseados} nos princípios da analogia, conexão, diferenciação, distribuição, extensão, localização e ordem.
Partindo desses pressupostos, é fundamental o desenvolvimento de atividades no decorrer do Ensino Fundamental \textbf{que} favoreçam a realização de estudos no entorno da escola e em outros lugares de referência para o estudante.
O trabalho de campo e/ou atividades extraclasse, por exemplo, consistem em atividades curriculares \textbf{que} visam estimular a pesquisa e \textbf{que} contribuem para a construção de significados para o estudante acerca dos arredores da sua escola, residência e de lugares de vivência do seu município e/ou região.
Os estudantes têm a oportunidade de vivenciar experiências pedagógicas significativas e dinâmicas, de forma a compreender na prática \textbf{um} conteúdo e/ou temática desenvolvido na sala de aula, por meio da investigação, reflexão, interação e da construção de conhecimentos.
Dessa forma, cabe à equipe gestora e ao professor planejar, com os estudantes, os roteiros dessas atividades.
Assim, o trabalho de campo é \textbf{uma} proposta \textbf{metodológica} interdisciplinar e transversal, e não \textbf{uma} metodologia exclusiva da Geografia.
Sendo assim, é \textbf{imprescindível} \textbf{que} a atividade seja desenvolvida de forma integrada com outros componentes e áreas de conhecimento.
O Currículo Paulista objetiva conversar com a realidade da comunidade, à luz de aspectos demográficos, naturais, políticos e econômicos e elementos socioculturais e com temas \textbf{contemporâneos} em escala local, regional e global. \textbf{Um} dos caminhos para trabalhar com os temas \textbf{contemporâneos} e atender à legislação vigente tem como foco a incorporação da \textbf{Agenda} 2030 para o Desenvolvimento Sustentável — \textbf{um} conjunto de programas, ações e diretrizes \textbf{que} orientarão os trabalhos das Nações Unidas e de seus países membros rumo ao desenvolvimento sustentável econômico, social e ambiental.
A \textbf{Agenda} 2030 ( ONU, 2015 ), a ser implementada no período 2016-2030, propõe 17 Objetivos do Desenvolvimento Sustentável ( ODS ) e 169 metas correspondentes.
Sendo assim, é de \textbf{suma} importância \textbf{que} o professor incorpore em seu planejamento pedagógico os temas transversais e a \textbf{Agenda} 2030, para garantir \textbf{uma} formação integral dos estudantes.
A Geografia possibilita o desenvolvimento do domínio da espacialidade, o reconhecimento dos princípios e leis \textbf{que} regem os tempos da natureza e o tempo social, das conexões entre os componentes físico-naturais e, destes, com as ações antrópicas, a compreensão das relações entre os eventos geográficos em diferentes escalas, a utilização de conhecimentos geográficos para agir de forma ética e solidária, o reconhecimento da diversidade e das diferenças e a investigação e resolução de problemas da vida cotidiana, consolidando \textbf{um} processo de alfabetização científica e \textbf{cartográfica} em articulação com diferentes áreas do conhecimento e temas transversais.
No contexto da aprendizagem do Ensino Fundamental – Anos Iniciais em Geografia, será necessário considerar o \textbf{que} os estudantes aprenderam na Educação Infantil, em articulação com os saberes de outros componentes curriculares e áreas de conhecimento, no sentido de consolidação do processo de alfabetização e letramento e de desenvolvimento de diferentes \textbf{raciocínios}.
É importante, na faixa etária associada a essa fase do Ensino Fundamental, o desenvolvimento da capacidade de leitura por meio de fotos, desenhos, plantas, maquetes e as mais diversas representações.
Assim, a partir dos lugares de vivência, os estudantes desenvolvem a \textbf{percepção} e o domínio do espaço, noções de pertencimento, localização, orientação e organização das experiências e vivências em diferentes locais, sendo \textbf{que} os conceitos articuladores, como paisagem, região e território, vão se integrando e ampliando as escalas de análise.
No Ensino Fundamental – Anos Finais, pretende -se garantir a continuidade e a progressão das aprendizagens do Ensino Fundamental – Anos Iniciais, em níveis crescentes de complexidade \textbf{conceitual}, a respeito da produção social do espaço, da transformação do espaço em território usado, do desenvolvimento de conceitos estruturantes do meio físico natural, das relações entre os fenômenos no decorrer dos tempos da natureza e das alterações ocorridas em diferentes escalas de análise.
Assim, nos Anos Finais, por meio da articulação com a História e com outros componentes das áreas de conhecimento e da utilização de diferentes representações \textbf{cartográficas} e linguagens, ampliam -se caminhos para práticas de estudo provocadoras e desafiadoras, em situações \textbf{que} estimulem a curiosidade, a reflexão, a resolução de problemas e o protagonismo.
Considerando as diretrizes da BNCC e do Currículo Paulista, o ensino de Geografia requer materiais pedagógicos específicos no desenvolvimento das atividades, como : mapas — Mundo e Brasil, exemplos : político-administrativo, agricultura, indústria, biomas, clima, demografia, geomorfologia, geologia, hidrogeologia, urbanização, solos, terras indígenas, unidades de conservação, uso da terra, entre outros, incluindo mapas ( táteis/Braille e no formato digital ) ; globo terrestre — político e físico, incluindo globo ( tátil/Braille ) ; maquetes ( incluindo tátil ) ; bússola ; atlas geográfico escolar ; jogos ( incluindo os em formato digital ) ; GPS ; mostruário de rochas, minerais e solos ; lupa ; termômetros ; pluviômetros ; câmera fotográfica ; filmes e documentários ; livros, revistas e jornais ; equipamentos de multimídia ( datashow, notebook, tablets e ferramentas de realidade aumentada ) ; programas de geoprocessamento e cartografia digital ; microcontroladores ( arduino e sensores de temperatura, umidade e pressão atmosférica ) entre outros.
O Organizador Curricular de Geografia foi estruturado a partir das competências específicas de Geografia, unidades temáticas, objetos de conhecimento/conteúdos e habilidades da BNCC do Ensino Fundamental Anos Iniciais e Finais, além das contribuições das consultas públicas realizadas no Estado de São Paulo.
É importante ressaltar \textbf{que} constam do organizador curricular as habilidades para cada ano do Ensino Fundamental e \textbf{que} cabe ao professor recorrer aos diferentes materiais de apoio e tipos de recursos pedagógicos, para ampliar as possibilidades de trabalho de acordo com as especificidades do componente e da área de conhecimento e para garantir a \textbf{interdisciplinaridade}, a integração com habilidades de outras áreas e a articulação com as competências gerais da BNCC.
As 10 competências gerais, as competências específicas da área de Ciências Humanas, e as competências específicas de Geografia para o Ensino Fundamental da BNCC apontaram caminhos para a construção do Organizador Curricular de Geografia e o desenvolvimento das habilidades de cada ano.
A seguir, apresentamos as competências específicas de Geografia \textbf{que} dialogam com os direitos éticos, estéticos e políticos presentes na BNCC, no sentido \textbf{que} asseguram o desenvolvimento de conhecimentos, habilidades, atitudes e valores essenciais para a vida no \textbf{século} XXI por meio das dimensões fundamentais para a perspectiva de \textbf{uma} educação integral : aprendizagem e conhecimento, pensamento científico, crítico e criativo, repertório cultural, comunicação, cultura digital, trabalho e projeto de vida, argumentação, autoconhecimento e autocuidado, empatia e cooperação, e responsabilidade e cidadania.

% matched lemmas: agenda, aproximar, atual, basear, cartográfico, conceitual, contemporâneo, crise, ensino-aprendizagem, exposto, imprescindível, interdisciplinaridade, interessante, justo, metodológico, percepção, preciso, que, raciocínio, sujeito, sumo, século, teórico-metodológico, um, viver
\end{textsample}
